% Part: computability
% Chapter: recursive-functions
% Section: bounded-minimization

\documentclass[../../../include/open-logic-section]{subfiles}

\begin{document}

\olfileid{cmp}{rec}{bmi}
\olsection{परिबद्ध लघुतमीकरण}

\begin{explain}
एखादा गुणधर्म किंवा संबंध~$P$ पूर्ण करणारी लघुतम संख्या म्हणून फलनाची
व्याख्या करणे अनेकदा उपयोगी ठरते. $P$ निर्णेय असेल, तर $P$ पूर्ण करणारी
लघुतम संख्या सापडेपर्यंत $0$, $1$, $2$, \dots, या सर्व संभाव्य संख्यांची
चाचणी घेऊन आपण हे फलन सहज संगणित करू शकतो. या प्रकारच्या अपरिबद्ध
शोधामुळे आपण आदिम पुनरावर्ती फलनांच्या क्षेत्राबाहेर जातो. मात्र स्वतंत्रपणे
दिलेल्या एखाद्या परिबंधापेक्षा \emph{लहान असलेल्या} लघुतम संख्येतच रस
असेल, तर आपण आदिम पुनरावर्ती क्षेत्रात राहतो. दुसऱ्या शब्दांत, आणि थोडे
अधिक सर्वसाधारणपणे, आपल्याकडे आदिम पुनरावर्ती संबंध~$R(x,z)$ आहे असे
समजा. $x$ आणि~$y$ यांना $R(x, z)$ सत्य करणाऱ्या लघुतम $z < y$ कडे
पाठवणारे फलन विचारात घ्या. $R(x, 0)$, $R(x, 1)$, \dots,
$R(x, y-1)$ यांची चाचणी घेऊन तेही संगणित करता येते. पण ते आदिम
पुनरावर्ती का आहे?
\end{explain}

\begin{prop}
$R(\vec x, z)$ आदिम पुनरावर्ती असेल, तर $R(\vec x, z)$ सत्य करणारे
लघुतम~$z$, $y$ पेक्षा लहान असा साक्षीदार असल्यास, परत करणारे आणि अन्यथा
$y$ परत करणारे फलन $m_R(\vec{x}, y)$ देखील आदिम पुनरावर्ती असते.
फलन~$m_R$ आपण असे लिहू:
\[
\bmin{z < y}{R(\vec{x}, z)},
\]
\end{prop}

\begin{proof}
$R(\vec{x}, z)$ सत्य करणारे~$z < 0$ असू शकत नाही, कारण मुळातच कोणतेही
$z < 0$ नसते. म्हणून $m_R(\vec x, 0) = 0$.

परिबंध $y + 1$ या रूपाचा असेल, तर तीन प्रकरणे असतात:
\begin{enumerate}
\item $R(\vec{x}, z)$ सत्य करणारे काही $z < y$ असते; अशा वेळी
$m_R(\vec{x}, y+1) = m_R(\vec{x}, y)$.
\item असे कोणतेही~$z<y$ नसते, पण $R(\vec{x}, y)$ सत्य असते; अशा वेळी
$m_R(\vec{x}, y+1) = y$.
\item $R(\vec{x}, z)$ सत्य करणारे कोणतेही $z < y+1$ नसते; अशा वेळी
$m_R(\vec{x}, y+1) = y+1$.
\end{enumerate}
म्हणून आदिम पुनरावर्तनाने $m_R(\vec x, 0)$ ची व्याख्या पुढीलप्रमाणे
देता येते:
\begin{align*}
m_R(\vec{x}, 0) & = 0\\
m_R(\vec{x}, y+1) & =
\begin{cases}
m_R(\vec{x}, y) & \text{जर $m_R(\vec x, y) \neq y$ असेल}\\
y & \text{जर $m_R(\vec x, y) = y$ आणि $R(\vec{x}, y)$ असेल}\\
y+1 & \text{अन्यथा.}
\end{cases}
\end{align*}
$R(\vec x, z)$ सत्य करणारे काही $z<y$ असते तेव्हा आणि केवळ तेव्हाच
$m_R(\vec x, y) \neq y$ असते, हे लक्षात घ्या.
\end{proof}

\begin{prob}
$R(\vec x, z)$ आदिम पुनरावर्ती आहे असे समजा. $R(\vec{x}, z)$ सत्य
करणारे लघुतम~$z$, $y$ पेक्षा लहान असा साक्षीदार असल्यास, परत करणारे आणि
अन्यथा $0$ परत करणारे फलन $m'_R(\vec{x}, y)$ हे~$\Char{R}$ पासून
आदिम पुनरावर्तनाने परिभाषित करा.
\end{prob}

\end{document}
